\section{Linked Lists}

A \textbf{singly linked list} is a linear data structure made of \textbf{nodes}.  
Each node stores:
\begin{itemize}
  \item a value (data)
  \item a pointer/reference to the next node
\end{itemize}

The last node points to \texttt{null}.

\begin{figure}[h]
\centering
\begin{tikzpicture}[
  >=Stealth,
  pointer/.style={->, thick}
]

% --- Node drawing helper: (x,y), title, value, next ---
\newcommand{\ListNode}[5]{%
  % Outer box
  \draw[thick] (#1,#2) rectangle ++(4,2);

  % Header divider
  \draw[thick] (#1,#2+1.3) -- ++(4,0);

  % Body vertical divider (value | next)
  \draw[thick] (#1+2.4,#2) -- ++(0,1.3);

  % Header text
  \node at (#1+2,#2+1.65) {\textbf{#3}};

  % Value / Next text
  \node at (#1+1.2,#2+0.65) {#4};
  \node at (#1+3.2,#2+0.65) {#5};
}

% --- Draw nodes ---
\ListNode{0}{0}{Node1}{red}{}
\ListNode{5.2}{0}{Node2}{blue}{}
\ListNode{10.4}{0}{Node3}{green}{null}

% --- Straight arrows (from next-field center to next node) ---
\draw[pointer] (4,0.65) -- (5.2,0.65);
\draw[pointer] (9.2,0.65) -- (10.4,0.65);

\end{tikzpicture}
\caption{A singly linked list where each node points to the next node}
\end{figure}



\subsection{Memory Layout}
Unlike arrays, linked lists are \textbf{not stored contiguously in memory}.  
Nodes can be located anywhere in memory and are connected using pointers.

\subsection{Structure}
\begin{itemize}
  \item \textbf{Head}: first node in the list
  \item \textbf{Tail}: last node in the list (points to \texttt{null})
\end{itemize}


\subsection{Traversal}
To iterate through a linked list:
\begin{verbatim}
curr = head
while (curr != null):
    curr = curr.next
\end{verbatim}
Traversal takes $O(n)$ time, where $n$ is the number of nodes.

\subsection{Insertion and Deletion}

\textbf{Inserting a node (given previous node)}
\begin{verbatim}
    newNode.next = prev.next
    prev.next = newNode
\end{verbatim}
\textbf{Deleting a node (given previous node):}
\begin{verbatim}
    prev.next = prev.next.next
\end{verbatim}
\textbf{Inserting before the head:}
\begin{verbatim}
    newNode.next = head
    head = newNode
\end{verbatim}
\textbf{Deleting the head:}
\begin{verbatim}
    head = head.next
\end{verbatim}

Inserting or deleting a node itself is $O(1)$, but finding the previous node takes $O(n)$ due to the lack of random access, so the overall operation is $O(n)$ if a search is required.

\subsection{Time Complexity}

\begin{center}
\begin{tabular}{l c}
\toprule
\textbf{Operation} & \textbf{Time Complexity} \\
\midrule
Access by index & $O(n)$ \\
Insert at head & $O(1)$ \\
Insert at tail & $O(1)$ (with tail pointer) \\
Insert in middle & $O(n)$ \\
Delete at head & $O(1)$ \\
Delete at tail & $O(n)$ \\
Delete in middle & $O(n)$ \\
\bottomrule
\end{tabular}
\end{center}

\subsection{Doubly Linked Lists}
In a \textbf{doubly linked list}, each node stores two pointers:
\begin{itemize}
  \item \texttt{next}: points to the next node
  \item \texttt{prev}: points to the previous node
\end{itemize}

This allows traversals forwards or backwards and enables insertions and deletions to be performed without needing to find the previous node, resulting in $O(1)$ when deleting at the end. Insertions/Deletions in the middle of the list are still $O(n)$, as we still need to access the i-th element.